% Data Mining - Chapter 3: Data Regression
\documentclass[10pt]{article}
% ============================================================
% _preamble.tex -- Shared preamble for all DM chapter docs
% Usage: % ============================================================
% _preamble.tex -- Shared preamble for all DM chapter docs
% Usage: % ============================================================
% _preamble.tex -- Shared preamble for all DM chapter docs
% Usage: \input{../_preamble} AFTER \documentclass, BEFORE \begin{document}
% ============================================================

\usepackage{pslatex}                    % Times New Roman
\usepackage[
    paperheight=22.86cm,
    paperwidth=15.24cm,
    left=1.27cm, right=1.27cm,
    top=1.6cm,   bottom=1.6cm,
    headheight=1in
]{geometry}
\usepackage[utf8]{inputenc}
\usepackage[vietnamese]{babel}
\usepackage[hidelinks]{hyperref}
\usepackage{graphicx}
\usepackage{enumerate}
\usepackage{float}
\usepackage{longtable}
\usepackage{array}
\usepackage{setspace}
\usepackage{sectsty}
\usepackage[normalem]{ulem}
\usepackage{caption}
\usepackage{titlesec}
\usepackage{amsmath}
\usepackage{amssymb}
\usepackage{enumitem}

% ----- Typography settings -----
\sectionfont{\fontsize{12pt}{12pt}\selectfont}
\captionsetup{font={normalsize,rm}, labelfont=bf, skip=10pt}
\titlelabel{\thetitle.\quad}
\urlstyle{same}

% ----- Paragraph indent -----
\usepackage{indentfirst}        % indent the first paragraph of every section
\setlength{\parindent}{1.5em}   % indent width (~4 spaces at 10pt)
\setlength{\parskip}{0pt}       % no extra space between paragraphs
 AFTER \documentclass, BEFORE \begin{document}
% ============================================================

\usepackage{pslatex}                    % Times New Roman
\usepackage[
    paperheight=22.86cm,
    paperwidth=15.24cm,
    left=1.27cm, right=1.27cm,
    top=1.6cm,   bottom=1.6cm,
    headheight=1in
]{geometry}
\usepackage[utf8]{inputenc}
\usepackage[vietnamese]{babel}
\usepackage[hidelinks]{hyperref}
\usepackage{graphicx}
\usepackage{enumerate}
\usepackage{float}
\usepackage{longtable}
\usepackage{array}
\usepackage{setspace}
\usepackage{sectsty}
\usepackage[normalem]{ulem}
\usepackage{caption}
\usepackage{titlesec}
\usepackage{amsmath}
\usepackage{amssymb}
\usepackage{enumitem}

% ----- Typography settings -----
\sectionfont{\fontsize{12pt}{12pt}\selectfont}
\captionsetup{font={normalsize,rm}, labelfont=bf, skip=10pt}
\titlelabel{\thetitle.\quad}
\urlstyle{same}

% ----- Paragraph indent -----
\usepackage{indentfirst}        % indent the first paragraph of every section
\setlength{\parindent}{1.5em}   % indent width (~4 spaces at 10pt)
\setlength{\parskip}{0pt}       % no extra space between paragraphs
 AFTER \documentclass, BEFORE \begin{document}
% ============================================================

\usepackage{pslatex}                    % Times New Roman
\usepackage[
    paperheight=22.86cm,
    paperwidth=15.24cm,
    left=1.27cm, right=1.27cm,
    top=1.6cm,   bottom=1.6cm,
    headheight=1in
]{geometry}
\usepackage[utf8]{inputenc}
\usepackage[vietnamese]{babel}
\usepackage[hidelinks]{hyperref}
\usepackage{graphicx}
\usepackage{enumerate}
\usepackage{float}
\usepackage{longtable}
\usepackage{array}
\usepackage{setspace}
\usepackage{sectsty}
\usepackage[normalem]{ulem}
\usepackage{caption}
\usepackage{titlesec}
\usepackage{amsmath}
\usepackage{amssymb}
\usepackage{enumitem}

% ----- Typography settings -----
\sectionfont{\fontsize{12pt}{12pt}\selectfont}
\captionsetup{font={normalsize,rm}, labelfont=bf, skip=10pt}
\titlelabel{\thetitle.\quad}
\urlstyle{same}

% ----- Paragraph indent -----
\usepackage{indentfirst}        % indent the first paragraph of every section
\setlength{\parindent}{1.5em}   % indent width (~4 spaces at 10pt)
\setlength{\parskip}{0pt}       % no extra space between paragraphs


\begin{document}
\onehalfspacing
\renewcommand{\arraystretch}{1.5}

\newcommand{\ChapterNum}{3}
\newcommand{\ChapterTitle}{HỒI QUY DỮ LIỆU (DATA REGRESSION)}
% ============================================================
% _title.tex -- Shared title block for all DM chapter docs
% Required macros (define BEFORE % ============================================================
% _title.tex -- Shared title block for all DM chapter docs
% Required macros (define BEFORE % ============================================================
% _title.tex -- Shared title block for all DM chapter docs
% Required macros (define BEFORE \input{../_title}):
%   \ChapterNum    e.g. {5}
%   \ChapterTitle  e.g. {PHÂN CỤM DỮ LIỆU (DATA CLUSTERING)}
% ============================================================

\begin{center}
    {\fontsize{14pt}{14pt}\selectfont
        \textbf{KHAI PHÁ DỮ LIỆU -- CHƯƠNG \ChapterNum}\par}
    \vspace{4pt}
    {\fontsize{13pt}{13pt}\selectfont
        \textbf{\ChapterTitle}\par}
    \vspace{6pt}
    {\fontsize{10pt}{10pt}\selectfont
        Tổng hợp kiến thức, câu hỏi tự luận và trắc nghiệm\par}
\end{center}
):
%   \ChapterNum    e.g. {5}
%   \ChapterTitle  e.g. {PHÂN CỤM DỮ LIỆU (DATA CLUSTERING)}
% ============================================================

\begin{center}
    {\fontsize{14pt}{14pt}\selectfont
        \textbf{KHAI PHÁ DỮ LIỆU -- CHƯƠNG \ChapterNum}\par}
    \vspace{4pt}
    {\fontsize{13pt}{13pt}\selectfont
        \textbf{\ChapterTitle}\par}
    \vspace{6pt}
    {\fontsize{10pt}{10pt}\selectfont
        Tổng hợp kiến thức, câu hỏi tự luận và trắc nghiệm\par}
\end{center}
):
%   \ChapterNum    e.g. {5}
%   \ChapterTitle  e.g. {PHÂN CỤM DỮ LIỆU (DATA CLUSTERING)}
% ============================================================

\begin{center}
    {\fontsize{14pt}{14pt}\selectfont
        \textbf{KHAI PHÁ DỮ LIỆU -- CHƯƠNG \ChapterNum}\par}
    \vspace{4pt}
    {\fontsize{13pt}{13pt}\selectfont
        \textbf{\ChapterTitle}\par}
    \vspace{6pt}
    {\fontsize{10pt}{10pt}\selectfont
        Tổng hợp kiến thức, câu hỏi tự luận và trắc nghiệm\par}
\end{center}


% ============================================================
\section{GIỚI THIỆU VỀ HỒI QUY}
% ============================================================

\subsection{Định nghĩa Hồi Quy}

\begin{itemize}
    \item \textbf{J. Han et al.:} Hồi quy là cơ chế thống kê cho phép dự đoán các giá trị số thực và liên tục.
    \item \textbf{R.D. Snee (1977):} Hồi quy là cơ chế thống kê cho phép dự đoán, kiểm soát và học các quy tắc sinh ra dữ liệu từ thực nghiệm.
    \item \textbf{Wiki (2009):} Phân tích hồi quy là cơ chế thống kê ước lượng mối quan hệ giữa các biến độc lập.
\end{itemize}

\textbf{Phân biệt quan trọng:}
\begin{itemize}
    \item \textbf{Hồi quy (Regression):} Dự đoán giá trị số thực (real-valued output) -- ví dụ: giá nhà, giá cổ phiếu.
    \item \textbf{Phân lớp (Classification):} ``Dự đoán'' cho giá trị rời rạc (discrete values) -- ví dụ: spam/not spam.
\end{itemize}

\subsection{Mô hình hồi quy}

Phương trình hồi quy tổng quát:
\[
Y = f(X, \theta)
\]
\begin{itemize}
    \item $X$: tập biến độc lập/dự báo (predictors/independent variables) -- mô tả các thay đổi của $Y$.
    \item $Y$: biến đáp ứng/phụ thuộc (response/dependent variable) -- mô tả sự kiện quan tâm.
    \item $\theta$: hệ số hồi quy (regression coefficients) -- mô tả mức độ ảnh hưởng tương đối của $X$ lên $Y$.
\end{itemize}

\subsection{Phân loại mô hình hồi quy}

\begin{itemize}
    \item \textbf{Tuyến tính (Linear) vs. Phi tuyến (Non-linear):} Linear: quan hệ tuyến tính giữa các tham số ảnh hưởng lên $Y$; Non-linear: ngược lại.
    \item \textbf{Đơn biến (Univariate) vs. Đa biến (Multivariate):} $X = (X_1)$ vs. $X = (X_1, X_2, \ldots, X_k)$.
    \item \textbf{Tham số (Parametric), phi tham số (Nonparametric), bán tham số (Semiparametric):}
    \begin{itemize}
        \item Parametric: $Y = \theta_0 + \theta_1 X$ (số tham số hữu hạn).
        \item Nonparametric: $Y = \theta_0 + f(X)$ (số tham số vô hạn).
        \item Semiparametric: $Y = \theta_0 + \theta_1 X_1 + f(X_2)$ (hữu hạn tham số quan tâm).
    \end{itemize}
    \item \textbf{Đối xứng (Symmetric) vs. Bất đối xứng (Asymmetric):}
    \begin{itemize}
        \item Symmetric: mô hình mô tả (log-linear models).
        \item Asymmetric: mô hình dự đoán (generalized linear models).
    \end{itemize}
\end{itemize}

% ============================================================
\section{HỒI QUY TUYẾN TÍNH ĐƠN BIẾN}
% ============================================================

\subsection{Giả thuyết và hàm mục tiêu}

Mô hình: $h_\theta(x) = \theta_0 + \theta_1 x$

\textbf{Ký hiệu:}
\begin{itemize}
    \item $N$: số mẫu huấn luyện.
    \item $x^{(i)}, y^{(i)}$: giá trị đầu vào/đầu ra của mẫu thứ $i$.
\end{itemize}

\textbf{Sai số dự đoán (residual):} $e^{(i)} = h_\theta(x^{(i)}) - y^{(i)}$

\textbf{Hàm chi phí (Cost function -- MSE):}
\[
J(\theta_0, \theta_1) = \frac{1}{2N} \sum_{i=1}^{N} \left(h_\theta(x^{(i)}) - y^{(i)}\right)^2
\]

\textbf{Mục tiêu:} Tìm $(\theta_0^*, \theta_1^*)$ để tối thiểu hóa $J(\theta_0, \theta_1)$.

\subsection{Phương pháp Gradient Descent}

\textbf{Ý tưởng:} Xuất phát từ điểm ngẫu nhiên $(\theta_0, \theta_1)$, lặp lại cập nhật để giảm $J$:

\textbf{Thuật toán (cập nhật đồng thời $\theta_0$ và $\theta_1$):}
\[
\theta_j := \theta_j - \alpha \frac{\partial J}{\partial \theta_j}
\]
\[
\frac{\partial J}{\partial \theta_0} = \frac{1}{N}\sum_{i=1}^{N}(h_\theta(x^{(i)}) - y^{(i)}), \quad
\frac{\partial J}{\partial \theta_1} = \frac{1}{N}\sum_{i=1}^{N}(h_\theta(x^{(i)}) - y^{(i)}) \cdot x^{(i)}
\]

\textbf{Learning rate $\alpha$:}
\begin{itemize}
    \item $\alpha$ quá nhỏ: học chậm, mất nhiều vòng lặp.
    \item $\alpha$ quá lớn: khó hội tụ, $J$ có thể không giảm tại mỗi bước.
    \item Thực nghiệm: thử $\alpha \in \{0.001, 0.003, 0.01, 0.03, 0.1, 0.3, \ldots\}$.
\end{itemize}

\subsection{Công thức trực tiếp (Closed-form)}

Khi $\theta_0 = 0$: Dừng lại ở $\hat{\theta}_1 = \frac{\sum(x^{(i)} - \bar{x})(y^{(i)} - \bar{y})}{\sum(x^{(i)} - \bar{x})^2}$ và $\hat{\theta}_0 = \bar{y} - \hat{\theta}_1 \bar{x}$.

% ============================================================
\section{HỒI QUY TUYẾN TÍNH ĐA BIẾN}
% ============================================================

\subsection{Mô hình đa biến}

\[
h_\theta(x) = \theta_0 + \theta_1 x_1 + \theta_2 x_2 + \ldots + \theta_n x_n = \theta^T x
\]
(với $x_0 = 1$)

\textbf{Hàm chi phí:}
\[
J(\theta) = \frac{1}{2N}\sum_{i=1}^{N}(h_\theta(x^{(i)}) - y^{(i)})^2
\]

\textbf{Gradient descent đa biến:}
\[
\theta_j := \theta_j - \alpha \frac{1}{N}\sum_{i=1}^{N}(h_\theta(x^{(i)}) - y^{(i)}) \cdot x_j^{(i)}, \quad j = 0, 1, \ldots, n
\]

\subsection{Feature Scaling (Chuẩn hóa đặc trưng)}

\textbf{Vấn đề:} Các thuộc tính có thang đo khác nhau (ví dụ: kích thước nhà 0--2000 ft², số phòng 1--5) làm gradient descent hội tụ chậm do contour plot $J$ bị kéo dài.

\textbf{Giải pháp:} Chuẩn hóa tất cả đặc trưng về cùng thang đo, ví dụ $x_j' \approx [-1, 1]$:
\[
x_j' = \frac{x_j - \mu_j}{\text{range}_j} \quad \text{(mean normalization)}
\]

\subsection{Phương pháp Normal Equation}

Thay vì dùng gradient descent lặp, có thể tính trực tiếp:
\[
\theta = (X^T X)^{-1} X^T Y
\]
trong đó $X$ là ma trận $N \times (n+1)$ (thêm cột $x_0 = 1$), $Y$ là vector $N \times 1$.

\textbf{So sánh Gradient Descent và Normal Equation:}

\begin{singlespace}
\begin{longtable}{|p{0.35\textwidth}|p{0.47\textwidth}|}
\hline
\textbf{Gradient Descent} & \textbf{Normal Equation} \\
\hline
\endhead
Cần chọn $\alpha$ & Không cần chọn $\alpha$ \\
\hline
Cần nhiều vòng lặp & Không cần lặp \\
\hline
Hoạt động tốt khi $n$ lớn ($n = 10^6$) & Phải tính $(X^TX)^{-1}$ -- $O(n^3)$ \\
\hline
& Không hiệu quả khi $n > 10^4$ \\
\hline
\end{longtable}
\end{singlespace}

\textbf{Vấn đề không khả nghịch (Non-invertible):} Xảy ra khi: (1) các biến phụ thuộc tuyến tính (ví dụ: kích thước m và feet); (2) $n > N$ (nhiều biến hơn mẫu).

% ============================================================
\section{HỒI QUY PHI TUYẾN}
% ============================================================

Khi $Y$ là hàm phi tuyến theo tham số $\theta$:
\[
Y = f(X, \theta) \quad \text{(ví dụ: hàm mũ, logarit, Gaussian)}
\]

Xác định $\theta$ tối ưu bằng: thuật toán tối ưu cục bộ (local optimization) hoặc tối ưu toàn cục (global optimization) tối thiểu hóa tổng bình phương sai số.

% ============================================================
\section{ỨNG DỤNG VÀ VẤN ĐỀ TRONG HỒI QUY}
% ============================================================

\subsection{Ứng dụng}

\begin{itemize}
    \item \textbf{Tiền xử lý DM:} Làm mịn dữ liệu, loại bỏ nhiễu.
    \item \textbf{Nhiệm vụ DM:} Dự đoán giá trị số, phân tích mô tả.
    \item \textbf{Các lĩnh vực:} Sinh học, nông nghiệp, kinh tế, kinh doanh, tài chính, bảo hiểm, e-commerce, khoa học, robot, điều khiển tự động.
\end{itemize}

\subsection{Đánh giá mô hình hồi quy}

\paragraph{Phương pháp tách dữ liệu:}
\begin{itemize}
    \item \textbf{Holdout method:} Chia ngẫu nhiên $D$ thành tập huấn luyện (2/3) và tập kiểm tra (1/3).
    \item \textbf{K-fold cross-validation:} Chia $D$ thành $k$ phần bằng nhau. Lặp $k$ lần: dùng phần $k$ để kiểm tra, phần còn lại để huấn luyện. Tính trung bình độ chính xác.
\end{itemize}

\paragraph{Các độ đo đánh giá độ chính xác:}
\begin{itemize}
    \item \textbf{SSE (Sum of Squared Errors):} $\text{SSE} = \sum_{i=1}^{n}(\hat{y}_i - y_i)^2$ -- đo lường tổng thể, nhỏ hơn tốt hơn.
    \item \textbf{MSE (Mean Squared Error):} $\text{MSE} = \frac{\text{SSE}}{n-m}$ ($m$: số hệ số hồi quy) -- đo phương sai còn lại, nhỏ hơn tốt hơn.
    \item \textbf{S (Standard Error of Estimate):} $S = \sqrt{\text{MSE}}$ -- sai số trung bình trong dự đoán, đo độ chính xác của dự đoán.
\end{itemize}

\subsection{Vấn đề trong hồi quy}

\begin{itemize}
    \item Giả định về phân phối dữ liệu (quan hệ giữa biến độc lập và phụ thuộc).
    \item Tính độc lập của các biến dự báo.
    \item Biến phải liên tục (cả predictors và responses).
    \item Lượng dữ liệu không đủ lớn.
    \item Xác định dạng mô hình hồi quy phù hợp.
    \item \textbf{Kỹ thuật nâng cao:} ANN (Artificial Neural Network), SVM (Support Vector Machine).
\end{itemize}

\paragraph{Các yếu tố ảnh hưởng đến thành công của mô hình hồi quy:}
\begin{enumerate}
    \item Xác định bài toán đúng đắn.
    \item Chọn biến quan trọng và dạng mô hình.
    \item Tập dữ liệu tốt (cả về khối lượng và chất lượng).
    \item Thuật toán ước lượng hệ số tốt (ví dụ: gradient descent).
    \item Kỹ thuật kiểm định mô hình.
\end{enumerate}

% ============================================================
\section{TÓM TẮT}
% ============================================================

\begin{itemize}
    \item Hồi quy là kỹ thuật thống kê dự đoán giá trị số thực liên tục.
    \item Phân loại: Linear/Non-linear, Univariate/Multivariate, 
    
    Parametric/Nonparametric/Semiparametric.
    \item Đơn biến: $h_\theta(x) = \theta_0 + \theta_1 x$; tối thiểu hóa $J$ bằng gradient descent.
    \item Đa biến: $h_\theta(x) = \theta^T x$; cần feature scaling; có thể dùng normal equation.
    \item Đánh giá: SSE, MSE, S; holdout hoặc k-fold cross-validation.
    \item Đơn giản nhưng hữu ích; ứng dụng rộng rãi; là ví dụ điển hình về đóng góp của thống kê vào DM.
\end{itemize}

% ============================================================
\newpage
\section{CÂU HỎI TỰ LUẬN}
% ============================================================

\begin{enumerate}[leftmargin=*, label=\textbf{Câu \arabic*.}]

    \item Hồi quy (Regression) là gì? Phân biệt hồi quy với phân lớp (Classification). Nêu ít nhất 3 ví dụ ứng dụng thực tiễn của hồi quy trong khai phá dữ liệu.

    \item Giải thích phương trình hồi quy $Y = f(X, \theta)$. Các thành phần $X$, $Y$, $\theta$ biểu diễn điều gì? Ý nghĩa của dấu và độ lớn của hệ số $\theta_i$ là gì?

    \item Trình bày đầy đủ các loại mô hình hồi quy: linear/nonlinear, univariate/multivariate, parametric/nonparametric/semiparametric. Cho ví dụ cụ thể cho mỗi loại.

    \item Hàm chi phí (Cost function) $J(\theta_0, \theta_1)$ trong hồi quy tuyến tính đơn biến là gì? Tại sao sử dụng MSE thay vì MAE (Mean Absolute Error)?

    \item Mô tả chi tiết thuật toán \textbf{Gradient Descent} cho hồi quy đơn biến. Tại sao cần cập nhật $\theta_0$ và $\theta_1$ \textbf{đồng thời}? Điều gì xảy ra nếu cập nhật tuần tự?

    \item Learning rate $\alpha$ trong Gradient Descent ảnh hưởng đến quá trình học như thế nào? Trình bày các trường hợp $\alpha$ quá nhỏ, $\alpha$ quá lớn và $\alpha$ phù hợp. Cách chọn $\alpha$ thực tế?

    \item Trình bày mô hình hồi quy tuyến tính đa biến. Sự khác biệt về thuật toán Gradient Descent giữa đơn biến và đa biến là gì?

    \item \textbf{Feature Scaling} là gì? Tại sao Feature Scaling quan trọng trong hồi quy đa biến với Gradient Descent? Mô tả hiệu ứng của thiếu scaling lên contour plot của $J$.

    \item So sánh \textbf{Gradient Descent} và \textbf{Normal Equation} để tối ưu hồi quy đa biến. Khi nào nên dùng mỗi phương pháp?

    \item Giải thích vấn đề \textbf{ma trận không khả nghịch (Non-invertible)} trong Normal Equation. Nguyên nhân và cách khắc phục như thế nào?

    \item Phương trình Normal Equation là gì? Dẫn xuất công thức $\theta = (X^TX)^{-1}X^TY$ và giải thích ý nghĩa của từng ma trận.

    \item Hồi quy phi tuyến (Non-linear Regression) khác hồi quy tuyến tính như thế nào? Cho 3 ví dụ về hàm phi tuyến và tình huống thực tế phù hợp.

    \item Trình bày phương pháp \textbf{K-fold Cross-Validation} để đánh giá mô hình hồi quy. Ưu điểm của K-fold so với Holdout method là gì? Giá trị $k$ nào phổ biến?

    \item Giải thích SSE, MSE và S (Standard Error of Estimate). Mối quan hệ giữa 3 độ đo này là gì? Khi nào thì mô hình được coi là có độ chính xác tốt?

    \item Trình bày ví dụ hồi quy đa biến với tập dữ liệu bán hàng: $y = \text{Quantity Sold} = 8536.21 - 835.72 \times \text{Price} + 0.592 \times \text{Advertising}$. Giải thích ý nghĩa của các hệ số.

    \item Giải thích tại sao hồi quy tuyến tính đơn biến với dạng $h_\theta(x) = \theta_0 + \theta_1 x$ tạo ra \textbf{convex cost function} J và không có local minima?

    \item Trình bày ứng dụng của hồi quy trong \textbf{tiền xử lý dữ liệu} (cụ thể trong Data Smoothing và Noise Removal). Hồi quy được dùng như thế nào để loại bỏ nhiễu?

    \item Phân biệt \textbf{Symmetric regression} và \textbf{Asymmetric regression}. Loại nào phù hợp cho nhiệm vụ dự đoán trong DM? Cho ví dụ cụ thể.

    \item Nêu 5 yếu tố ảnh hưởng đến sự thành công của mô hình hồi quy. Yếu tố nào bạn cho là quan trọng nhất và tại sao?

    \item So sánh hồi quy tuyến tính truyền thống với \textbf{ANN (Artificial Neural Network)} và \textbf{SVM (Support Vector Machine)} trong bối cảnh hồi quy nâng cao. Khi nào cần dùng các phương pháp nâng cao?

\end{enumerate}

% ============================================================
\newpage
\section{CÂU HỎI TRẮC NGHIỆM}
% ============================================================

\begin{enumerate}[leftmargin=*, label=\textbf{Câu \arabic*.}]

    \item Hồi quy (Regression) trong DM dùng để dự đoán:
    \begin{enumerate}[label=\Alph*.]
        \item Giá trị rời rạc (nhãn lớp).
        \item Giá trị số thực và liên tục.
        \item Các tập phổ biến trong tập giao dịch.
        \item Số lượng cụm trong tập dữ liệu.
    \end{enumerate}

    \item Mô hình hồi quy $Y = \theta_0 + \theta_1 X$ thuộc loại:
    \begin{enumerate}[label=\Alph*.]
        \item Nonparametric univariate.
        \item Parametric univariate linear.
        \item Semiparametric multivariate.
        \item Nonlinear parametric.
    \end{enumerate}

    \item Hàm chi phí $J(\theta_0, \theta_1)$ trong hồi quy tuyến tính được tối:
    \begin{enumerate}[label=\Alph*.]
        \item Tối đa hóa.
        \item Tối thiểu hóa.
        \item Giữ cố định.
        \item Tối đa hóa ở bước đầu, tối thiểu hóa ở bước sau.
    \end{enumerate}

    \item Sai số dự đoán (residual) $e^{(i)}$ được định nghĩa là:
    \begin{enumerate}[label=\Alph*.]
        \item $y^{(i)} - \bar{y}$
        \item $h_\theta(x^{(i)}) - y^{(i)}$
        \item $x^{(i)} - \bar{x}$
        \item $\theta_1 x^{(i)} - \theta_0$
    \end{enumerate}

    \item Trong Gradient Descent, learning rate $\alpha$ quá lớn dẫn đến:
    \begin{enumerate}[label=\Alph*.]
        \item Hội tụ rất nhanh đến nghiệm tối ưu.
        \item Hội tụ chậm, cần nhiều vòng lặp.
        \item $J(\theta)$ có thể không giảm tại mỗi bước, thuật toán có thể không hội tụ.
        \item Gradient descent luôn hội tụ bất kể $\alpha$.
    \end{enumerate}

    \item Trong Gradient Descent đơn biến, cập nhật $\theta_0$ và $\theta_1$ phải:
    \begin{enumerate}[label=\Alph*.]
        \item Cập nhật $\theta_0$ trước, sau đó mới dùng $\theta_0$ mới để cập nhật $\theta_1$.
        \item Cập nhật $\theta_1$ trước, sau đó cập nhật $\theta_0$.
        \item Cập nhật đồng thời (simultaneously update).
        \item Cập nhật luân phiên theo từng vòng lặp.
    \end{enumerate}

    \item Feature Scaling trong hồi quy đa biến giúp:
    \begin{enumerate}[label=\Alph*.]
        \item Tăng độ chính xác của mô hình hồi quy.
        \item Gradient Descent hội tụ nhanh hơn bằng cách cân bằng thang đo các đặc trưng.
        \item Giảm overfitting.
        \item Tự động chọn giá trị $\alpha$ tối ưu.
    \end{enumerate}

    \item Normal Equation $\theta = (X^TX)^{-1}X^TY$ có ưu điểm so với Gradient Descent là:
    \begin{enumerate}[label=\Alph*.]
        \item Hiệu quả hơn khi $n$ rất lớn ($n > 10^6$).
        \item Không cần chọn $\alpha$ và không cần vòng lặp.
        \item Luôn tìm được nghiệm tối ưu toàn cục dù $J$ không lồi.
        \item Có thể áp dụng khi ma trận $(X^TX)$ không khả nghịch.
    \end{enumerate}

    \item Ma trận $(X^TX)$ trong Normal Equation \textbf{không khả nghịch} khi:
    \begin{enumerate}[label=\Alph*.]
        \item Số mẫu $N$ quá nhỏ và $n < N$.
        \item Tồn tại biến phụ thuộc tuyến tính hoặc số biến $n$ lớn hơn số mẫu $N$.
        \item Tất cả biến đều độc lập tuyến tính.
        \item $\alpha$ được chọn quá nhỏ.
    \end{enumerate}

    \item Mô hình hồi quy $Y = \theta_0 + f(X)$ thuộc loại:
    \begin{enumerate}[label=\Alph*.]
        \item Parametric.
        \item Nonparametric.
        \item Semiparametric.
        \item Linear.
    \end{enumerate}

    \item Độ đo nào sau đây đo lường \textbf{tổng thể} sai số của mô hình hồi quy?
    \begin{enumerate}[label=\Alph*.]
        \item MSE.
        \item S (Standard Error of Estimate).
        \item SSE (Sum of Squared Errors).
        \item $R^2$.
    \end{enumerate}

    \item MSE (Mean Squared Error) được tính bằng:
    \begin{enumerate}[label=\Alph*.]
        \item $\text{SSE} \times (n - m)$
        \item $\text{SSE} / (n - m)$
        \item $\sqrt{\text{SSE}}$
        \item $\text{SSE} / n$
    \end{enumerate}

    \item Trong đánh giá mô hình hồi quy, \textbf{K-fold cross-validation} với $k=10$ có nghĩa là:
    \begin{enumerate}[label=\Alph*.]
        \item Chia dữ liệu thành 10 phần, mỗi lần dùng 9 phần test và 1 phần train.
        \item Chia dữ liệu thành 10 phần; mỗi lần dùng 1 phần test, 9 phần train; lặp 10 lần.
        \item Lặp gradient descent 10 lần.
        \item Chọn 10 giá trị $\alpha$ khác nhau và lấy trung bình.
    \end{enumerate}

    \item Hồi quy tuyến tính đơn biến dạng $h_\theta(x) = \theta_0 + \theta_1 x^2$ thuộc loại:
    \begin{enumerate}[label=\Alph*.]
        \item Phi tuyến về tham số (nonlinear in parameters).
        \item Tuyến tính về tham số (linear in parameters).
        \item Phi tham số (nonparametric).
        \item Bán tham số (semiparametric).
    \end{enumerate}

    \item Dấu của hệ số $\theta_i$ trong mô hình hồi quy biểu diễn:
    \begin{enumerate}[label=\Alph*.]
        \item Độ lớn tuyệt đối của ảnh hưởng $X_i$ lên $Y$.
        \item Hướng ảnh hưởng (tương quan thuận/nghịch) của $X_i$ lên $Y$.
        \item Mức độ tương quan giữa $X_i$ và các biến khác.
        \item Xác suất biến $X_i$ xuất hiện trong mô hình.
    \end{enumerate}

    \item Standard Error $S$ của mô hình hồi quy đo lường:
    \begin{enumerate}[label=\Alph*.]
        \item Tổng bình phương sai số.
        \item Sai số trung bình trong quá trình dự đoán, đo độ chính xác của dự đoán.
        \item Sai số khi ước lượng $\theta$ bằng gradient descent.
        \item Số vòng lặp cần thiết để hội tụ.
    \end{enumerate}

    \item Khi vẽ $J(\theta)$ theo số vòng lặp, thuật toán gradient descent \textbf{không hội tụ} khi:
    \begin{enumerate}[label=\Alph*.]
        \item $J(\theta)$ giảm đơn điệu qua mỗi vòng lặp.
        \item $J(\theta)$ giảm nhanh lúc đầu rồi chậm dần.
        \item $J(\theta)$ tăng hoặc dao động.
        \item $J(\theta)$ bằng 0.
    \end{enumerate}

    \item Với mô hình $y = 9323 - 823 \times \text{price}$ và $R^2 = 0.65$, $R^2$ cho biết:
    \begin{enumerate}[label=\Alph*.]
        \item $65\%$ biến động của $y$ được giải thích bởi mô hình.
        \item $65\%$ các dự đoán là chính xác.
        \item SSE bằng 0.65.
        \item Mô hình chính xác $65\%$ trường hợp.
    \end{enumerate}

    \item Phương pháp \textbf{Holdout} chia dữ liệu thành:
    \begin{enumerate}[label=\Alph*.]
        \item $k$ phần bằng nhau.
        \item Tập huấn luyện và tập kiểm tra (thường 2/3 và 1/3).
        \item Tập huấn luyện, tập validation và tập kiểm tra.
        \item 10 fold để cross-validation.
    \end{enumerate}

    \item Hồi quy tuyến tính đơn biến với $J(\theta_0, \theta_1) = \text{MSE}$ có hình dạng của $J$ là:
    \begin{enumerate}[label=\Alph*.]
        \item Hàm lõm (concave function).
        \item Hàm lồi (convex function) -- dạng bowl/paraboloid.
        \item Hàm có nhiều điểm cực tiểu cục bộ.
        \item Hàm tuyến tính phẳng.
    \end{enumerate}

    \item Với hồi quy đa biến, khi $n = 10^6$ (triệu biến), nên dùng:
    \begin{enumerate}[label=\Alph*.]
        \item Normal equation vì không cần vòng lặp.
        \item Gradient descent vì Normal equation tính $(X^TX)^{-1}$ rất tốn kém.
        \item Cả hai đều hiệu quả như nhau.
        \item Chỉ có thể dùng K-NN regression.
    \end{enumerate}

    \item Kỹ thuật hồi quy nâng cao nào được đề cập trong slide C3?
    \begin{enumerate}[label=\Alph*.]
        \item K-Means và DBSCAN.
        \item Logistic Regression và Decision Tree.
        \item ANN (Artificial Neural Network) và SVM (Support Vector Machine).
        \item Apriori và FP-Growth.
    \end{enumerate}

    \item Mô hình $Y = \theta_0 + \theta_1 X_1 + f(X_2)$ thuộc loại:
    \begin{enumerate}[label=\Alph*.]
        \item Parametric.
        \item Nonparametric.
        \item Semiparametric.
        \item Linear nonparametric.
    \end{enumerate}

    \item Trong hồi quy, nếu $\theta_1 = -835.72$ (hệ số của biến Price), điều đó có nghĩa:
    \begin{enumerate}[label=\Alph*.]
        \item Giá tăng 1 đơn vị thì Quantity Sold tăng 835.72.
        \item Giá tăng 1 đơn vị thì Quantity Sold giảm 835.72.
        \item Quantity Sold luôn bằng $-835.72 \times \text{Price}$.
        \item Price và Quantity Sold không có quan hệ tuyến tính.
    \end{enumerate}

    \item Hồi quy tuyến tính \textbf{đơn biến} có giả thuyết:
    \begin{enumerate}[label=\Alph*.]
        \item $h_\theta(x) = \theta_0 + \theta_1 x_1 + \theta_2 x_2$
        \item $h_\theta(x) = \theta_0 + \theta_1 x$
        \item $h_\theta(x) = g(\theta^T x)$
        \item $h_\theta(x) = \theta^T x + f(x)$
    \end{enumerate}

    \item Trong Gradient Descent, điều kiện dừng (convergence) thường là:
    \begin{enumerate}[label=\Alph*.]
        \item Số vòng lặp đạt đúng $n$.
        \item $J(\theta)$ không thay đổi đáng kể giữa các vòng lặp (hoặc thay đổi dưới ngưỡng).
        \item $\theta_0 = \theta_1 = 0$.
        \item $J(\theta) = 0$.
    \end{enumerate}

    \item Ứng dụng của hồi quy trong \textbf{tiền xử lý dữ liệu} là:
    \begin{enumerate}[label=\Alph*.]
        \item Phân lớp điểm dữ liệu vào các nhóm.
        \item Làm mịn và loại bỏ nhiễu dữ liệu.
        \item Tìm các luật kết hợp.
        \item Phân cụm dữ liệu.
    \end{enumerate}

    \item Khi dùng Normal Equation với $n = 4$ biến và $N = 4$ mẫu, rủi ro lớn nhất là:
    \begin{enumerate}[label=\Alph*.]
        \item Gradient descent không hội tụ.
        \item $(X^TX)$ có thể không khả nghịch do $n \approx N$ (quá ít mẫu so với biến).
        \item $\alpha$ quá lớn.
        \item Dữ liệu không có phân phối Gaussian.
    \end{enumerate}

    \item Tổng bình phương sai số (SSE) trong đánh giá hồi quy được tính bằng:
    \begin{enumerate}[label=\Alph*.]
        \item $\sum_{i=1}^{n} |\hat{y}_i - y_i|$
        \item $\sum_{i=1}^{n} (\hat{y}_i - y_i)^2$
        \item $\frac{1}{n}\sum_{i=1}^{n}(\hat{y}_i - y_i)^2$
        \item $\sqrt{\frac{\text{SSE}}{n-m}}$
    \end{enumerate}

    \item Mô hình $h_\theta(x) = \theta_0 + \theta_1 x + \theta_2 x^2$ biểu diễn:
    \begin{enumerate}[label=\Alph*.]
        \item Hồi quy phi tuyến về tham số.
        \item Hồi quy tuyến tính về tham số (linear in parameters) với $x^2$ là một đặc trưng.
        \item Hồi quy phi tham số.
        \item Không thể xác định.
    \end{enumerate}

    \item Trong cross-validation, mục đích của việc dùng dữ liệu không thấy trong quá trình huấn luyện (unseen test data) để đánh giá là:
    \begin{enumerate}[label=\Alph*.]
        \item Kiểm tra khả năng tổng quát hóa (generalization) của mô hình.
        \item Tối ưu hóa thêm các tham số $\theta$.
        \item Giảm SSE của tập huấn luyện.
        \item Tăng tốc độ huấn luyện.
    \end{enumerate}

\end{enumerate}

% ============================================================
\newpage
\section*{ĐÁP ÁN}
% ============================================================

\subsection*{Câu hỏi tự luận -- Hướng dẫn trả lời}

\begin{singlespace}
\begin{longtable}{|p{0.08\textwidth}|p{0.82\textwidth}|}
\hline
\textbf{Câu} & \textbf{Nội dung cần trình bày} \\
\hline
\endhead
1 & Hồi quy: dự đoán giá trị số thực. Khác phân lớp ở output (liên tục vs. rời rạc). Ví dụ: dự đoán giá nhà, giá cổ phiếu, doanh số bán hàng. \\
\hline
2 & $X$: biến độc lập/dự báo; $Y$: biến phụ thuộc/đáp ứng; $\theta$: hệ số. Dấu $\theta_i$: hướng ảnh hưởng (+: thuận, -: nghịch). Độ lớn $|\theta_i|$: mức độ ảnh hưởng. \\
\hline
3 & Linear: $Y = \theta_0 + \theta_1 X$; Nonlinear: $Y = e^{\theta_1 X}$; Univariate: 1 biến đầu vào; Multivariate: nhiều biến; Parametric: hữu hạn tham số; Nonparametric: $Y = f(X)$; Semiparametric: $Y = \theta_1 X + f(X_2)$. \\
\hline
4 & $J = \frac{1}{2N}\sum(h_\theta(x^{(i)}) - y^{(i)})^2$. MSE hơn MAE vì: khả vi liên tục (dễ lấy đạo hàm), phạt sai số lớn mạnh hơn, tối ưu hóa thuận lợi hơn. \\
\hline
5 & Khởi tạo $(\theta_0, \theta_1)$; lặp: tính gradient cho cả $\theta_0$ và $\theta_1$ bằng tham số hiện tại; cập nhật đồng thời. Nếu tuần tự: $\theta_1$ cập nhật sau sẽ dùng $\theta_0$ mới $\Rightarrow$ gradient tính sai, sai kết quả. \\
\hline
6 & $\alpha$ nhỏ: $J$ giảm chậm, nhiều vòng lặp. $\alpha$ lớn: $J$ tăng/dao động, không hội tụ. Chọn thực tế: thử 0.001, 0.003, 0.01, ... theo cấp số nhân; vẽ $J$ vs. số vòng lặp để kiểm tra. \\
\hline
7 & Đa biến: $h_\theta(x) = \theta^T x = \theta_0 + \sum \theta_j x_j$. Gradient descent cập nhật $n+1$ tham số; công thức gradient: $\frac{\partial J}{\partial \theta_j} = \frac{1}{N}\sum(h_\theta(x^{(i)}) - y^{(i)})x_j^{(i)}$. \\
\hline
8 & Feature Scaling: chuẩn hóa đặc trưng về $\approx [-1,1]$. Thiếu scaling: contour plot của $J$ bị kéo dài hình elip, gradient descent đi theo đường ziczac $\Rightarrow$ hội tụ chậm. \\
\hline
9 & Gradient descent: cần chọn $\alpha$, lặp nhiều; tốt với $n$ lớn. Normal equation: không cần $\alpha$, không lặp; phải tính $(X^TX)^{-1}$ $O(n^3)$; không dùng khi $n > 10^4$. \\
\hline
10 & Non-invertible khi: (1) biến phụ thuộc tuyến tính (ví dụ: kích thước m và feet$^2$) $\to$ xóa biến trùng; (2) $n > N$ $\to$ giảm số biến hoặc thu thập thêm dữ liệu. \\
\hline
11 & Normal equation: đặt gradient = 0: $\frac{\partial J}{\partial \theta} = 0$ $\Rightarrow$ $X^TX\theta = X^TY$ $\Rightarrow$ $\theta = (X^TX)^{-1}X^TY$. $X$: ma trận đặc trưng $N \times (n+1)$; $Y$: vector nhãn $N \times 1$; $(X^TX)^{-1}$: giả nghịch đảo. \\
\hline
12 & Phi tuyến: $Y$ là hàm phi tuyến theo $\theta$. Ví dụ: (1) $Y = \theta_0 e^{\theta_1 X}$ -- tăng trưởng; (2) $Y = \theta_0 \log(\theta_1 X)$ -- suy giảm logarit; (3) $Y = \frac{1}{1+e^{-\theta^TX}}$ -- sigmoid. \\
\hline
13 & K-fold: chia thành $k$ phần; lặp $k$ lần, mỗi lần dùng 1 phần test; tính trung bình độ chính xác. Ưu điểm so với holdout: ít phụ thuộc cách chia dữ liệu, sử dụng toàn bộ dữ liệu để train/test. $k=10$ phổ biến nhất. \\
\hline
14 & SSE: tổng bình phương sai số (overall measure). MSE = SSE/(n-m): phương sai còn lại (nhỏ = ít phương sai chưa giải thích). $S = \sqrt{\text{MSE}}$: sai số trung bình dự đoán (đơn vị giống $Y$). Mô hình tốt khi SSE, MSE, S nhỏ. \\
\hline
15 & $\theta_0 = 8536.21$: giá trị cơ sở khi Price=0 và Advertising=0 (không có ý nghĩa thực tế ở đây). $\theta_1 = -835.72$: Giá tăng 1\$ thì QS giảm 835.72 (quan hệ nghịch). $\theta_2 = 0.592$: Quảng cáo tăng 1\$ thì QS tăng 0.592. \\
\hline
16 & $J(\theta_0, \theta_1)$ là hàm toàn phương (quadratic) -- hình bowl (paraboloid lồi). Hàm lồi không có điểm cực tiểu cục bộ, chỉ có 1 cực tiểu toàn cục $\Rightarrow$ gradient descent luôn tìm được nghiệm tối ưu. \\
\hline
17 & Hồi quy trong preprocessing: (1) Data smoothing -- dùng đường hồi quy để làm mịn dữ liệu (thay giá trị thực bằng giá trị trên đường hồi quy); (2) Noise removal -- xác định điểm xa đường hồi quy là nhiễu và loại bỏ. \\
\hline
18 & Symmetric regression: mô hình mô tả -- không phân biệt biến phụ thuộc/độc lập (ví dụ: log-linear). Asymmetric: mô hình dự đoán -- phân biệt rõ $X$ và $Y$ (ví dụ: generalized linear). Trong DM dự đoán, dùng asymmetric. \\
\hline
19 & 5 yếu tố: (1) Xác định bài toán đúng đắn; (2) Chọn biến và dạng mô hình; (3) Dữ liệu tốt; (4) Thuật toán ước lượng tốt; (5) Kỹ thuật kiểm định. Quan trọng nhất: dữ liệu tốt -- garbage in, garbage out. \\
\hline
20 & Hồi quy tuyến tính: đơn giản, giải thích được, hoạt động tốt khi quan hệ tuyến tính. ANN: xử lý phi tuyến phức tạp, nhiều tham số, ít giải thích. SVM: hiệu quả với chiều cao, robust với outlier. Dùng nâng cao khi: quan hệ phi tuyến mạnh, dữ liệu phức tạp, yêu cầu độ chính xác cao. \\
\hline
\end{longtable}
\end{singlespace}

\subsection*{Câu hỏi trắc nghiệm -- Đáp án}

\begin{singlespace}
\begin{longtable}{|c|c||c|c||c|c||c|c|}
\hline
\textbf{Câu} & \textbf{ĐA} & \textbf{Câu} & \textbf{ĐA} & \textbf{Câu} & \textbf{ĐA} & \textbf{Câu} & \textbf{ĐA} \\
\hline
\endhead
1 & B & 11 & C & 21 & B & 31 & B \\
\hline
2 & B & 12 & B & 22 & C & 32 & B \\
\hline
3 & B & 13 & B & 23 & C & 33 & B \\
\hline
4 & B & 14 & B & 24 & B & 34 & B \\
\hline
5 & C & 15 & B & 25 & B & 35 & B \\
\hline
6 & C & 16 & B & 26 & B & 36 & B \\
\hline
7 & B & 17 & C & 27 & B & 37 & B \\
\hline
8 & B & 18 & A & 28 & C & 38 & A \\
\hline
9 & B & 19 & B & 29 & B & 39 & B \\
\hline
10 & B & 20 & B & 30 & B & 40 & A \\
\hline
\end{longtable}
\end{singlespace}

\end{document}
